% TOP_PROBLEM: {"id": 100, "title": "Roth's Theorem with Random Common Differences", "category_id": 2, "category": {"id": 2, "name": "combinatorics", "display_name": "Combinatorics", "description": "Counting problems, graph theory, discrete structures.", "slug": "combinatorics", "order_index": 2, "created_at": "2026-07-31T15:26:25.671Z"}, "status": "open"}
% FIRST_POSTED: null
% ENTRY_KIND: statement_only
% Imported from: batch02.tex; UnsolvedMath ID 100
\documentclass[11pt]{article}
\usepackage[margin=25mm]{geometry}
\usepackage{fontspec}
\setmainfont{Latin Modern Roman}
\usepackage{amsmath,amssymb,mathtools}
\usepackage{unicode-math}
\setmathfont{Latin Modern Math}
\usepackage{xurl,hyperref}
\hypersetup{colorlinks=true,urlcolor=blue}
\setcounter{secnumdepth}{0}
\setlength{\parindent}{0pt}
\setlength{\parskip}{5pt}
\setlength{\emergencystretch}{3em}
\newcommand{\sref}[2]{\hyperlink{source:#1:#2}{[S#2]}}
\newcommand{\cataloguescope}[1]{\paragraph{Recorded target.}#1}
\begin{document}
% UnsolvedMath ID 100; source catalogue code GREEN-010
\setcounter{section}{99}
\section{Roth's Theorem with Random Common Differences}\label{100}
{\small\sffamily UnsolvedMath ID 100 · Combinatorics}
\subsection{Definitions and mathematical statement}

\paragraph{Shared notation.}$\mathbb N=\{1,2,\ldots\}$, and $\mathbb Z,\mathbb Q,\mathbb R,\mathbb C$ denote the integers, rationals, reals and complex numbers. Any other conventions are specified locally.


Let $(p_n)_{n\ge1}$ be a sequence in $[0,1]$. Form a random subset $S\subset\mathbb N$ by including each $n$ independently with probability $p_n$. A set $S$ is a \emph{set of two-recurrence} if every $A\subset\mathbb N$ of positive upper density,
\[
\overline d(A)=\limsup_{N\to\infty}\frac{|A\cap\{1,\ldots,N\}|}{N}>0,
\]
contains three distinct terms $a,a+d,a+2d$ with $d\in S$. The probability-one assertion concerns this simultaneous property for every such $A$, rather than a separate event for each fixed $A$.
\paragraph{Statement recorded in the supplied source.}
Determine conditions on the inclusion probabilities $(p_n)$ under which $S$ is almost surely a set of two-recurrence. In particular, does this hold whenever $p_n=\min\{1,\omega(n)/n\}$ for a function $\omega:\mathbb N\to(0,\infty)$ satisfying $\omega(n)\to\infty$? The required conclusion is that, with probability one, every positive-upper-density set of integers contains a nontrivial three-term arithmetic progression whose common difference lies in $S$. \sref{100}{1}\sref{100}{2}
\subsection{Short English statement}
Which independent random sets of allowed positive common differences preserve Roth's theorem for every positive-density set? The optimistic target permits inclusion probabilities only a diverging factor above $1/n$.
\subsection{Sources}
\begin{description}
\item[\hypertarget{source:100:1}{S1}] ULAM.AI and the UnsolvedMath Contributors, \emph{UnsolvedMath: A Curated Collection of Open Mathematics Problems}, record 100 (GREEN-010). CC BY 4.0. \url{https://huggingface.co/datasets/ulamai/UnsolvedMath}.
\item[\hypertarget{source:100:2}{S2}] Ben Green, \emph{100 Open Problems}, author manuscript with December 2025 updates, Problem 10 and its comments, pp. 7--8. \url{https://people.maths.ox.ac.uk/greenbj/papers/open-problems.pdf}.
\end{description}
\label{end:100}
\end{document}
